% Source-aligned English rebound to repaired p.119 French.
% Decisions: TOH-P119-R1-EN-001 through -006.
\SourcePage{119}
\LayerTarget{tohoku:paper}
\begin{center}
  {\Large\bfseries ON SOME POINTS OF HOMOLOGICAL ALGEBRA.\par}
  \vspace{1.2em}
  {\large\scshape Alexandre Grothendieck\par}
  \vspace{1.2em}
  {(Received March 1, 1957)\par}
\end{center}

\LayerTarget{tohoku:introduction}
\LayerContents{Introduction}
\section*{Introduction\textsuperscript{1)}}

\LayerTarget{tohoku:introduction:I}
\LayerTarget{tohoku:introduction:I:p1}
\paragraph{I. Contents of the work.}
This work arose from an attempt to exploit the formal analogy between the
theory of the cohomology of a space with coefficients in a sheaf [4], [5] and
the theory of the derived functors of functors on modules [6], in order to find
a common framework in which these theories and others could be encompassed.

\LayerTarget{tohoku:introduction:I:p2}
This framework is outlined in Chapter I, whose theme is the same as that of
[3]. These two accounts, however, overlap only in no.~1.4. I have devoted
particular attention to giving workable criteria, using the notion of infinite
sums and products in abelian categories, for the existence of ``enough''
injective or projective objects in an abelian category, without which the
essential homological techniques cannot be applied. Moreover, for the
convenience of the reader, considerable space has been devoted to presenting
the language of functors (nos.~1.1, 1.2, and 1.3). \emph{The introduction of
additive categories} in no.~1.3, as a preliminary to abelian categories,
provides a convenient language (for example, for treating spectral functors in
Chapter II).

\LayerTarget{tohoku:introduction:I:p3}
Chapter II outlines the essential points of the homological formalism in
abelian categories. The publication of [6] allowed me to be very concise,
since the techniques of Cartan--Eilenberg carry over without any change to the
new framework. Nos.~2.1 and 2.2 were written, however, so as not to exclude
abelian categories that do not contain enough injective or projective objects.
In the following numbers we make full use of the customary techniques of
resolutions. Nos.~2.4 and 2.5 contain various supplements and are essential for
understanding what follows. In particular, Theorem 2.4.1 gives a mechanical
way to obtain most of the known spectral sequences (and in any event
\emph{all} those encountered in this work).

\LayerTarget{tohoku:introduction:I:p4}
In Chapter III we redevelop the theory of the cohomology of a space with
coefficients in a sheaf, including the classical Leray spectral sequences. The
account given here represents a relaxation compared with [4], [15], in
particular in that all the essential results are obtained without imposing, at
almost any point in this chapter (nor in the subsequent ones), a restrictive
hypothesis on the nature of the spaces under consideration; so that

\LayerTarget{tohoku:introduction:footnote:1}
\begin{quote}
\footnotesize
1) The essential content of Chapters I, II, IV, and part of Chapter III was
developed in the spring of 1955, on the occasion of a seminar on Homological
Algebra given at the University of Kansas. ---The numbers in brackets refer to
the end of this article.
\end{quote}

\TohokuContinuation{120}
